% ============================================================
% 10 -- ROOT CAUSE, FIXES, AND LESSONS LEARNED
% ============================================================

\section{Root Cause, Fixes, and Lessons Learned}
\label{sec:lessons}

LAMP was not compromised because of one spectacular bug.

The exploit worked because attacker-controlled data repeatedly crossed
boundaries between different interpreters without being constrained for
the language it entered next.

The complete chain involved:

\begin{center}
  \displayfont
  HTTP
  $\longrightarrow$
  TeX
  $\longrightarrow$
  SQL
  $\longrightarrow$
  TeX
  $\longrightarrow$
  shell
\end{center}

At each step, a value that looked like ordinary application data in one
context gained new meaning in another.


% ------------------------------------------------------------
% 10.1
% ------------------------------------------------------------

\subsection{Root Cause: Data Became Code}

The most important vulnerability was not
\code{--shell-escape} by itself.

It was the loss of separation between data and executable TeX.

The component type is the clearest example.

It began as a POST value:

\begin{lstlisting}
typ=light
\end{lstlisting}

It was escaped for SQL and persisted into MariaDB.

Later, the same value was reconstructed through:

\begin{lstlisting}
\edef\type#1{\compdata[3]}
\end{lstlisting}

At that point it was no longer merely database content.

It could influence TeX expansion.

Once attacker-controlled TeX existed inside a process started with shell
escape enabled, the impact became operating-system command execution.

\begin{findingbox}
Escaping a value for one interpreter does not make it safe for the next
interpreter.

The component type was SQL-safe but not TeX-safe.
\end{findingbox}


% ------------------------------------------------------------
% 10.2
% ------------------------------------------------------------

\subsection{Fixing the Path Collision}

The \code{Path} vulnerability existed because HTTP headers and internal
request variables shared the same dynamically constructed namespace.

A header named:

\begin{lstlisting}
Path:
\end{lstlisting}

became:

\begin{lstlisting}
\httpPath
\end{lstlisting}

which collided directly with application state.

A later defensive wrapper rejected header names corresponding to
internal request variables before the request reached XeLaTeX:

\begin{lstlisting}
path
method
version
query
static
line
cookies
\end{lstlisting}

That prevents the specific overwrite used by the exploit.

A cleaner design would go further and avoid placing external HTTP header
names and internal application state in the same namespace at all.


% ------------------------------------------------------------
% 10.3
% ------------------------------------------------------------

\subsection{Fixing the Component Injection}

The web interface only supported three component types:

\begin{lstlisting}
light
button
source
\end{lstlisting}

The vulnerable backend nevertheless accepted arbitrary values.

The effective fix was therefore not complicated:

\begin{lstlisting}
typ =~ /^(?:light|button|source)$/
\end{lstlisting}

This is preferable to trying to blacklist TeX syntax.

There are many ways to express interesting behaviour in TeX, and trying
to enumerate dangerous tokens would leave unnecessary parser complexity.

If a field has three valid values, then it should accept exactly three
valid values.


% ------------------------------------------------------------
% 10.4
% ------------------------------------------------------------

\subsection{Fixing the Newline Carrier}

The original \code{\textbackslash store} function attempted to escape
characters with special meaning in TeX:

\begin{lstlisting}
\
{
}
_
^
\end{lstlisting}

The exploit did not need any of them.

It only needed a newline.

The patched request handling instead restricted ship names to a small
set of expected characters before they reached the TeX application.

That removed the ability to transform:

\begin{lstlisting}
\def \shipname {VALUE}
\end{lstlisting}

into a multi-line file containing an independent shell command.

Again, the important lesson is that allowlisting the expected data is
simpler than attempting to escape every dangerous representation.


% ------------------------------------------------------------
% 10.5
% ------------------------------------------------------------

\subsection{Shell Escape Multiplied the Impact}

The TeX injection was the vulnerability.

Shell escape was what made its consequences severe.

LAMP was executed with:

\begin{lstlisting}
xelatex -8bit --shell-escape /srv/main.tex
\end{lstlisting}

Once arbitrary TeX expansion was achieved, the distance from TeX code
execution to operating-system command execution became very small.

If shell escape is not required, the obvious hardening measure is to
disable it.

If external command execution is genuinely required, the safer design
is to expose only the specific operations the application needs rather
than giving general TeX processing access to a shell.


% ------------------------------------------------------------
% 10.6
% ------------------------------------------------------------

\subsection{Defense in Depth}

The competition patches closed the primitives used by this exploit.

There are still several additional measures that would reduce the impact
of similar mistakes.

\begin{itemize}

  \item Run XeLaTeX as a dedicated unprivileged user rather than root.

  \item Give the process access only to files required by the service.

  \item Keep attacker-controlled values as inert data instead of
        reconstructing them as TeX source.

  \item Validate input before passing it into the TeX interpreter.

  \item Separate HTTP request metadata from internal TeX control
        sequences.

  \item Prefer strict allowlists for small enumerated fields such as
        component types.

  \item Disable unrestricted shell escape whenever possible.

\end{itemize}

None of these replaces correct input handling.

They reduce the damage when correct input handling eventually fails.


% ------------------------------------------------------------
% 10.7
% ------------------------------------------------------------

\subsection{The Exploit in One Page}

The entire chain can finally be summarized as:

\begin{center}
\begin{tabular}{c}
  \code{Path:} overwrites \code{\textbackslash httpPath} \\[0.25em]
  $\downarrow$ \\[0.25em]
  arbitrary file retrieval \\[0.25em]
  $\downarrow$ \\[0.25em]
  attacker-controlled \code{typ} stored in MariaDB \\[0.25em]
  $\downarrow$ \\[0.25em]
  database value reconstructed as TeX \\[0.25em]
  $\downarrow$ \\[0.25em]
  \code{\^{}\^{}5cUseName\{@@input\}} reaches shell escape \\[0.25em]
  $\downarrow$ \\[0.25em]
  container-root command execution \\[0.25em]
  $\downarrow$ \\[0.25em]
  newline-injected ship file becomes long-command carrier \\[0.25em]
  $\downarrow$ \\[0.25em]
  recent checker storage searched \\[0.25em]
  $\downarrow$ \\[0.25em]
  flags retrieved through the original \code{Path} primitive \\[0.25em]
  $\downarrow$ \\[0.25em]
  MOTH submission
\end{tabular}
\end{center}


% ------------------------------------------------------------
% 10.8
% ------------------------------------------------------------

\subsection{Final Takeaway}

LAMP is a good example of why unusual software stacks produce unusual
security boundaries.

HTTP headers became TeX macros.

Database strings became TeX macros.

User-controlled text became files that were later interpreted again.

And the TeX interpreter had access to a shell.

Individually, several of those decisions were manageable.

Together, they produced a chain from one HTTP request field to root
command execution inside the service container.

\begin{findingbox}
The useful question was never:

\begin{center}
``Is LaTeX dangerous?''
\end{center}

It was:

\begin{center}
``At what point does attacker-controlled data become executable LaTeX?''
\end{center}

Once that boundary was found, the rest of LAMP went brrr.
\end{findingbox}